% ADD THIS HEADER TO ALL NEW CHAPTER FILES FOR SUBFILES SUPPORT

% Allow independent compilation of this section for efficiency
\documentclass[../CLthesis.tex]{subfiles}

% Add the graphics path for subfiles support
\graphicspath{{\subfix{../images/}}}

% END OF SUBFILES HEADER

%%%%%%%%%%%%%%%%%%%%%%%%%%%%%%%%%%%%%%%%%%%%%%%%%%%%%%%%%%%%%%%%
% START OF DOCUMENT: Every chapter can be compiled separately
\begin{document}

\chapter{Experiment I: BERTopic-Based Discourse Classification}%
\label{chap:bertopic}

This chapter presents the first experiment in our analytical framework: an unsupervised topic modelling of the multilingual discussion corpus using BERTopic. The goal is to identify the major thematic structures present across English, Japanese, and Chinese posts, without imposing any predefined category scheme. The resulting topic map provides a foundation for the subsequent supervised analyses in Chapters~\ref{chap:who}--\ref{chap:why}, which investigate who participates in these discourses, how they communicate, and why specific positions are adopted.

BERTopic was applied to the full multilingual corpus. The model first embeds each document using a multilingual sentence transformer (\texttt{[MODEL PLACEHOLDER]}), then reduces dimensionality with UMAP before clustering via HDBSCAN. Topic representations are finally extracted using a class-based TF-IDF weighting scheme. The key hyperparameters were set as follows: UMAP \texttt{n\_neighbors = [X]}, \texttt{n\_components = [X]}, \texttt{min\_dist = [X]}; HDBSCAN \texttt{min\_cluster\_size = [X]}, \texttt{min\_samples = [X]}. A total of 47 topics were extracted, with 2{,}530 documents assigned to the outlier class (Topic~$-1$). Figure~\ref{fig:datamap} provides an overview of the resulting topic landscape, where each point represents a document and colours indicate topic membership.

\begin{figure}[htbp]
    \centering
    % \includegraphics[width=\textwidth]{RQ1/datamap_bertopic_labelled_cluster.png}
    \caption{Document map of the BERTopic classification results. Each point represents one document; colours correspond to distinct topics. Topic labels are manually assigned based on keyword inspection.}
    \label{fig:datamap}
\end{figure}

The 47 topics vary considerably in size (ranging from 20 to 499 documents) and in language composition. Rather than discussing every topic in detail, the following sections highlight a curated selection: three topics dominated by a single language community, three cross-lingual topics that draw participants from all three communities, and an analysis of the inter-topic relationships revealed by hierarchical clustering and pairwise similarity.

%---------------------------------------------------------------
\section{Monolingual and Cross-Lingual Topics}
\label{sec:mono-cross}

To characterise the thematic landscape, we distinguish between \textit{monolingual topics} --- those in which one language community accounts for the overwhelming majority of documents --- and \textit{cross-lingual topics}, where all three communities contribute meaningfully. Figure~\ref{fig:langdist} shows the language distribution across all 47 topics, making this distinction visually apparent.

\begin{figure}[htbp]
    \centering
    % \includegraphics[width=\textwidth]{RQ1/language_distribution_per_topic.png}
    \caption{Language composition per topic (stacked bar chart). Blue = English, orange = Japanese, green = Chinese. Topics with one dominant colour are monolingual; those with mixed colours are cross-lingual.}
    \label{fig:langdist}
\end{figure}

\subsection{Monolingual Topics}

Three topics stand out as essentially monolingual, each dominated by a distinct language community. They are selected as the topic with the highest share of each respective language across all 47 clusters.

\textbf{Topic~0 --- Unification Church and Japanese Politics (Japanese, 98.8\%).}
Topic~0 is by far the largest topic in the corpus (499 documents) and is almost exclusively Japanese (entropy $\approx 0.07$). Its defining keywords are 統一教会 (\textit{Unification Church}), 安倍 (\textit{Abe}), 日本 (\textit{Japan}), 韓国 (\textit{Korea}), and 自民党 (\textit{LDP}). Representative documents revolve around the political entanglement between the Japanese Liberal Democratic Party --- and particularly former Prime Minister Shinzo Abe --- and the Holy Spirit Association for the Unification of World Christianity. One post argues that Abe, having been fully aware of the church's inner workings, bore direct responsibility as its public promoter. Another traces how Abe's faction within the LDP allowed the church to rename itself, thereby enabling it to evade a prosecution blacklist and continue its ``spirit sales'' fraud against Japanese believers. A third calls for believers to withhold financial donations and turn their backs on clergy during sermons, framing the LDP–church nexus as a source of social corruption. The near-total absence of English or Chinese contributions confirms that this controversy resonates almost exclusively within the Japanese public sphere, shaped by the shock of the 2022 assassination and the ensuing media scrutiny of decades-long ties between conservative politicians and the movement.

\textbf{Topic~16 --- Religion, Immigration, and Canadian Identity (English, 91.8\%).}
Topic~16 (61 documents; entropy $\approx 0.34$) is the most English-dominant topic in the dataset. Key terms include \textit{canada}, \textit{party}, \textit{immigrants}, \textit{country}, \textit{canadian}, \textit{conservatives}, and \textit{society}. One post carefully distinguishes between the orderly, family-oriented immigration system in Canada and the chaotic mass migration seen elsewhere in the West, pushing back against conflation of the two. A second post reflects on historical barriers faced by Canadians before the 1960s, questioning whether the country's identity should remain tethered to a ``British and Christian'' self-image. A third post makes this tension explicit, asserting that Canada was built on European Christian civilisational foundations and that immigrants are expected to integrate into that tradition. Together, these documents show how English-speaking discussants situate religion within a broader debate over national identity, multiculturalism, and conservative political values --- a framing that has little resonance in the Japanese or Chinese communities.

\textbf{Topic~43 --- Missionaries and the Spread of Christianity in China (Chinese, 95.5\%).}
Topic~43 (22 documents; entropy $\approx 0.18$) is the most Chinese-dominant topic in the corpus. Its keywords --- 传教士 (\textit{missionaries}), 传教 (\textit{evangelism}), 中国 (\textit{China}), 道教 (\textit{Taoism}), 儒教 (\textit{Confucianism}), 洋人 (\textit{foreigners}) --- centre on the historical and comparative question of how religions propagate themselves in China. Representative documents contrast the institutional missionary infrastructure of Christianity (both Catholic and Protestant) with Chinese traditional religions: one post notes that Confucianism did not need missionaries because it was promoted by imperial decree from the Han dynasty onward; another argues pithily that the imperial examination system itself functioned as the greatest ``missionary'' for Confucian values; a third draws a systematic comparison across Christianity, Buddhism, and Taoism to highlight that the absence of a missionary tradition in Confucianism and Taoism reflects their fundamentally different logics of propagation. This topic is almost entirely absent from English and Japanese threads, reflecting a distinctly Chinese intellectual concern with the historical interaction between indigenous belief systems and Western Christian missions.

\subsection{Cross-Lingual Topics}

Three topics display the highest language entropy in the corpus, meaning that all three communities contribute in roughly equal measure. These are selected on the basis of entropy score and represent thematic concerns that cut across linguistic and cultural boundaries.

\textbf{Topic~36 --- Conservative vs.\ Liberal Factions within the Church (entropy = 1.097).}
Topic~36 (25 documents) is the most evenly distributed topic in the entire dataset: Chinese (36\%), English (32\%), and Japanese (32\%). Its keywords --- 保守派 (\textit{conservatives}), 政治 (\textit{politics}), 投票 (\textit{voting}), 左派 (\textit{leftists}), 骑士团/騎士団 (\textit{Knights}) --- converge on the internal political fracture lines of mainstream Christianity. Representative documents from all three languages articulate a strikingly similar diagnosis: progressive and mainline churches are losing ground on both ends of the political spectrum. A Chinese post observes that mainstream denominations are being squeezed out --- attacked from the right for being too liberal and abandoned by secular leftists who have no interest in faith at all --- and that this leaves them commanding only around 17\% of the American religious landscape. A Japanese post, almost word-for-word parallel in its argument, describes the same dynamic using Pew Research figures and laments that liberal Christians are now ignored by both the conservative Catholic--Evangelical alliance and the secular progressive movement. An English post engages with a related debate about whether conservative rhetoric in church disputes mirrors left-wing censorship, while also questioning the sincerity of those who deploy grand Catholic chivalric identities (Knights of the Holy Sepulchre, Knights of Malta) while showing little actual faith commitment. The cross-lingual symmetry of this topic suggests that the conservative--liberal fault line within Christianity is experienced as a genuinely global institutional tension, not a culturally specific one.

\textbf{Topic~3 --- LGBTQ+ Identity and Catholic Teaching (entropy = 1.095).}
Topic~3 (149 documents) has a near-perfectly balanced distribution: Chinese (36\%), English (34\%), and Japanese (30\%). Keywords include 同性恋/同性愛 (\textit{homosexuality}), \textit{gay}, \textit{lgbt}, \textit{disordered}, and 性 (\textit{sexuality}). Documents across all three languages engage critically with the Church's characterisation of same-sex attraction as ``intrinsically disordered.'' A Chinese post describes the rule barring openly gay men from the priesthood as forcing candidates into a ``pathological and dangerous self-hatred,'' and demands that bishops clarify whether they actually intend to enforce the policy --- or admit that it is merely performative. A second Chinese document takes a more theological angle, arguing that as long as discussion remains confined to a heterosexual frame, the Church avoids confronting fundamental questions about whether any sexual behaviour can be described as consistent with human dignity. A Japanese post attacks the natural law reasoning behind the prohibition as a rhetorical convenience --- ``a phrase used to make the opponent go silent without any actual argument'' --- and dismisses clerical sexual conservatism as reflecting personal anxiety rather than principled theology. The consistency of this critical discourse across three language communities indicates that LGBTQ+ inclusion is experienced as a universal fault line within Catholicism rather than a locally specific controversy.

\textbf{Topic~7 --- Science, Faith, and Creationism (entropy = 1.092).}
Topic~7 (119 documents) is similarly balanced: English (39\%), Japanese (30\%), and Chinese (31\%). Keywords include 科学/\textit{science}, 進化/进化 (\textit{evolution}), 創造/创造 (\textit{creation}), 理解 (\textit{understanding}), and \textit{evidence}. The topic brings together posts that debate the relationship between religious belief and scientific knowledge, with a particular focus on creationism and evolutionary theory. A Japanese post directly challenges an interlocutor who claims that ``original sin is encoded in DNA,'' calling this ``absurd'' and laying out the scientific consensus that human genetic diversity points to multiple ancestral populations, not a single couple. A Chinese post offers a sharper structural critique: creationists, it argues, do not merely hold different beliefs but reject the very definition of science accepted by the scientific community, treating the Bible as a literal instruction manual such that if it said $2+2=5$, they would be obliged to accept it and declare mathematics a lie. A second Japanese post grounds the argument historically, noting that Newton himself distinguished between science and religion, and that even the Pope has stated that religion is not the basis of science. What unites documents across all three communities is not a rejection of faith per se but a sustained critique of a literalist epistemology that refuses to engage with empirical evidence. The cross-lingual coherence of this topic suggests that the science--religion debate, at least in its anti-creationist form, is a shared concern among educated online Catholic discussants worldwide.

%---------------------------------------------------------------
\section{Hierarchical Relationships Between Topics}
\label{sec:hierarchy}

Figure~\ref{fig:hierarchy} presents the hierarchical clustering of the 47 topics, computed on the basis of topic embedding similarity. The dendrogram reveals which topics are semantically proximate and how thematic clusters are nested within broader super-topics.

\begin{figure}[htbp]
    \centering
    % \includegraphics[width=\textwidth]{RQ1/hierarchical_clustering_bertopic.png}
    \caption{Hierarchical clustering of BERTopic topics based on topic embedding similarity. Topics that merge at low distances (left side of the dendrogram) are closely related; those merging at greater distances share only broad thematic affinity.}
    \label{fig:hierarchy}
\end{figure}

Several groupings are visible in the dendrogram. At the top of the hierarchy, a broad cluster unites many of the Japanese-dominant topics --- including Topics 0 (Japan Church/State), 17, 31, and 23 --- reflecting the shared vocabulary and discourse norms of the Japanese-language threads. These topics merge at a relatively short distance, indicating that their linguistic surface features create a strong intra-group coherence even when their substantive content differs (e.g., church-state relations versus intra-Christian identity debates).

A second major cluster groups the ethically contested topics together: Topics~2 (Women's Ordination), 3 (LGBTQ+), 6 (Clerical Abuse), 12 (Divorce and Remarriage), and 27 (Gay and Law) tend to merge before connecting to politically oriented topics such as Topic~1 (Trump/Politics) or Topic~16 (Canada Immigration). This grouping suggests that the model correctly captures a distinction between \textit{intra-ecclesial moral debates} and \textit{religion-in-politics} discussions, even without being given any explicit category labels.

A third cluster of Chinese-dominant topics --- including Topics~20, 26, 37, and 43 --- also emerges as a coherent group at the level of language, paralleling the Japanese-language cluster. However, within this Chinese-language grouping, the dendrogram further separates topics by content: folk religion and Buddhism (Topics 37, 43) diverge earlier from discussions of Catholic-specific practices (Topics 44, 28).

Topics related to science and religion (Topic~7) and salvation doctrine (Topic~19) form a small but distinct cluster, merging with Catholic practice topics (Topic~5, liturgy/mass) only at a high hierarchical distance. This indicates that epistemological debates --- about evolution, evidence, and the compatibility of faith with reason --- occupy a relatively separate discursive space from liturgical or governance concerns within the corpus.

Overall, the hierarchical structure suggests that language of origin is one of the strongest organising forces in the data, followed by thematic domain (ethics, politics, practice, theology). This ordering will be important when interpreting the supervised classification results in later chapters.

%---------------------------------------------------------------
\section{Inter-Topic Similarity}
\label{sec:similarity}

Figure~\ref{fig:similarity} displays the pairwise cosine similarity matrix computed over topic embeddings. Each cell represents the semantic proximity between two topics; darker cells indicate higher similarity.

\begin{figure}[htbp]
    \centering
    % \includegraphics[width=0.85\textwidth]{RQ1/유사도_Ber.png}
    \caption{Pairwise cosine similarity matrix of topic embeddings. Diagonal entries (self-similarity = 1.0) are shown in dark blue; off-diagonal cells indicate the degree to which two topics share embedding space.}
    \label{fig:similarity}
\end{figure}

The matrix reveals that the vast majority of topic pairs have very low similarity scores, clustering near zero. This is consistent with BERTopic's HDBSCAN-based clustering mechanism, which pushes topic embeddings apart in the reduced dimensional space: once topics are formed, they tend to be well-separated by construction. The near-zero off-diagonal values confirm that the 47 topics are largely \textit{semantically distinct}, with little bleed-over between thematic clusters.

A small number of topic pairs show marginally elevated similarity. Notably, topics within the same language community --- particularly the Japanese-dominant and Chinese-dominant clusters identified in the hierarchical analysis --- tend to exhibit slightly higher pairwise similarity than cross-language topic pairs. This is consistent with the finding from Section~\ref{sec:hierarchy} that language acts as a strong organising factor. Beyond this language-induced similarity, however, no pair of substantively different topics reaches a cosine score high enough to suggest genuine thematic overlap.

In practical terms, the low inter-topic similarity implies that the 47 topics can be treated as largely independent units for downstream analysis: documents assigned to Topic~2 (Women's Ordination) and Topic~6 (Clerical Abuse), for example, share little latent content despite both falling under the broad umbrella of Catholic institutional critique. This granularity is a strength for our subsequent research questions, as it allows us to target specific discourse types with precision.

%---------------------------------------------------------------
\section{Summary and Transition}
\label{sec:bertopic-summary}

This chapter has established the thematic map of multilingual Catholic discourse through unsupervised BERTopic modelling. The 47 extracted topics range from highly language-specific clusters --- most notably the Japanese-dominant Topic~0 on the Unification Church–Abe controversy --- to genuinely cross-lingual debates on gender, sexuality, and institutional abuse that attract roughly equal participation from English, Japanese, and Chinese speakers.

Several structural insights emerge from this first experiment. First, \textit{language is the strongest organising factor} in the corpus: monolingual topics arise not only because of culturally specific events (e.g., the Abe assassination and its political fallout) but also because of distinct discourse norms and intra-community references that do not translate across linguistic boundaries. Second, \textit{certain ethical concerns transcend language}: the debates on women's ordination, LGBTQ+ inclusion, and clerical abuse converge across all three communities with strikingly similar critical frames, suggesting that these are experienced as universal institutional failings rather than culturally localised grievances. Third, \textit{the topics are semantically well-separated}, as confirmed by the low inter-topic similarity scores, which gives us confidence that subsequent analyses of specific topic clusters will not be confounded by category overlap.

These findings motivate the structure of the remaining chapters. Rather than treating the corpus as a single undifferentiated whole, we use the topic classifications established here to stratify the data for deeper investigation. Chapters~\ref{chap:who} through~\ref{chap:why} follow a \textit{who~$\rightarrow$ how~$\rightarrow$ why} logic: we first ask who participates in these discourses (participant profiles and community structures), then how they communicate (rhetorical and linguistic strategies), and finally why they take the positions they do (underlying value systems and argumentation patterns). The BERTopic topic labels serve as a principled organisational backbone for all three stages of this analysis.

% self-defined macro: include bibliography even when compiling a single chapter
\subfilebibliography
\end{document}
